\IfFileExists{stacks-project.cls}{%
\documentclass{stacks-project}
}{%
\documentclass{amsart}
}

% For dealing with references we use the comment environment
\usepackage{verbatim}
\newenvironment{reference}{\comment}{\endcomment}
%\newenvironment{reference}{}{}
\newenvironment{slogan}{\comment}{\endcomment}
\newenvironment{history}{\comment}{\endcomment}

% For commutative diagrams we use Xy-pic
\usepackage[all]{xy}

% We use 2cell for 2-commutative diagrams.
\xyoption{2cell}
\UseAllTwocells

% We use multicol for the list of chapters between chapters
\usepackage{multicol}

% This is generally recommended for better output
\usepackage{lmodern}
\usepackage[T1]{fontenc}

% For cross-file-references
\usepackage{xr-hyper}

% Package for hypertext links:
\usepackage{hyperref}

% For any local file, say "hello.tex" you want to link to please
% use \externaldocument[hello-]{hello}
\externaldocument[introduction-]{introduction}
\externaldocument[conventions-]{conventions}
\externaldocument[sets-]{sets}
\externaldocument[categories-]{categories}
\externaldocument[topology-]{topology}
\externaldocument[sheaves-]{sheaves}
\externaldocument[sites-]{sites}
\externaldocument[stacks-]{stacks}
\externaldocument[fields-]{fields}
\externaldocument[algebra-]{algebra}
\externaldocument[brauer-]{brauer}
\externaldocument[homology-]{homology}
\externaldocument[derived-]{derived}
\externaldocument[simplicial-]{simplicial}
\externaldocument[more-algebra-]{more-algebra}
\externaldocument[smoothing-]{smoothing}
\externaldocument[modules-]{modules}
\externaldocument[sites-modules-]{sites-modules}
\externaldocument[injectives-]{injectives}
\externaldocument[cohomology-]{cohomology}
\externaldocument[sites-cohomology-]{sites-cohomology}
\externaldocument[dga-]{dga}
\externaldocument[dpa-]{dpa}
\externaldocument[sdga-]{sdga}
\externaldocument[hypercovering-]{hypercovering}
\externaldocument[schemes-]{schemes}
\externaldocument[constructions-]{constructions}
\externaldocument[properties-]{properties}
\externaldocument[morphisms-]{morphisms}
\externaldocument[coherent-]{coherent}
\externaldocument[divisors-]{divisors}
\externaldocument[limits-]{limits}
\externaldocument[varieties-]{varieties}
\externaldocument[topologies-]{topologies}
\externaldocument[descent-]{descent}
\externaldocument[perfect-]{perfect}
\externaldocument[more-morphisms-]{more-morphisms}
\externaldocument[flat-]{flat}
\externaldocument[groupoids-]{groupoids}
\externaldocument[more-groupoids-]{more-groupoids}
\externaldocument[etale-]{etale}
\externaldocument[chow-]{chow}
\externaldocument[intersection-]{intersection}
\externaldocument[pic-]{pic}
\externaldocument[weil-]{weil}
\externaldocument[adequate-]{adequate}
\externaldocument[dualizing-]{dualizing}
\externaldocument[duality-]{duality}
\externaldocument[discriminant-]{discriminant}
\externaldocument[derham-]{derham}
\externaldocument[local-cohomology-]{local-cohomology}
\externaldocument[algebraization-]{algebraization}
\externaldocument[curves-]{curves}
\externaldocument[resolve-]{resolve}
\externaldocument[models-]{models}
\externaldocument[functors-]{functors}
\externaldocument[equiv-]{equiv}
\externaldocument[pione-]{pione}
\externaldocument[etale-cohomology-]{etale-cohomology}
\externaldocument[proetale-]{proetale}
\externaldocument[relative-cycles-]{relative-cycles}
\externaldocument[more-etale-]{more-etale}
\externaldocument[trace-]{trace}
\externaldocument[crystalline-]{crystalline}
\externaldocument[spaces-]{spaces}
\externaldocument[spaces-properties-]{spaces-properties}
\externaldocument[spaces-morphisms-]{spaces-morphisms}
\externaldocument[decent-spaces-]{decent-spaces}
\externaldocument[spaces-cohomology-]{spaces-cohomology}
\externaldocument[spaces-limits-]{spaces-limits}
\externaldocument[spaces-divisors-]{spaces-divisors}
\externaldocument[spaces-over-fields-]{spaces-over-fields}
\externaldocument[spaces-topologies-]{spaces-topologies}
\externaldocument[spaces-descent-]{spaces-descent}
\externaldocument[spaces-perfect-]{spaces-perfect}
\externaldocument[spaces-more-morphisms-]{spaces-more-morphisms}
\externaldocument[spaces-flat-]{spaces-flat}
\externaldocument[spaces-groupoids-]{spaces-groupoids}
\externaldocument[spaces-more-groupoids-]{spaces-more-groupoids}
\externaldocument[bootstrap-]{bootstrap}
\externaldocument[spaces-pushouts-]{spaces-pushouts}
\externaldocument[spaces-chow-]{spaces-chow}
\externaldocument[groupoids-quotients-]{groupoids-quotients}
\externaldocument[spaces-more-cohomology-]{spaces-more-cohomology}
\externaldocument[spaces-simplicial-]{spaces-simplicial}
\externaldocument[spaces-duality-]{spaces-duality}
\externaldocument[formal-spaces-]{formal-spaces}
\externaldocument[restricted-]{restricted}
\externaldocument[spaces-resolve-]{spaces-resolve}
\externaldocument[formal-defos-]{formal-defos}
\externaldocument[defos-]{defos}
\externaldocument[cotangent-]{cotangent}
\externaldocument[examples-defos-]{examples-defos}
\externaldocument[algebraic-]{algebraic}
\externaldocument[examples-stacks-]{examples-stacks}
\externaldocument[stacks-sheaves-]{stacks-sheaves}
\externaldocument[criteria-]{criteria}
\externaldocument[artin-]{artin}
\externaldocument[quot-]{quot}
\externaldocument[stacks-properties-]{stacks-properties}
\externaldocument[stacks-morphisms-]{stacks-morphisms}
\externaldocument[stacks-limits-]{stacks-limits}
\externaldocument[stacks-cohomology-]{stacks-cohomology}
\externaldocument[stacks-perfect-]{stacks-perfect}
\externaldocument[stacks-introduction-]{stacks-introduction}
\externaldocument[stacks-more-morphisms-]{stacks-more-morphisms}
\externaldocument[stacks-geometry-]{stacks-geometry}
\externaldocument[moduli-]{moduli}
\externaldocument[moduli-curves-]{moduli-curves}
\externaldocument[examples-]{examples}
\externaldocument[exercises-]{exercises}
\externaldocument[guide-]{guide}
\externaldocument[desirables-]{desirables}
\externaldocument[coding-]{coding}
\externaldocument[obsolete-]{obsolete}
\externaldocument[fdl-]{fdl}
\externaldocument[index-]{index}

% Theorem environments.
%
\theoremstyle{plain}
\newtheorem{theorem}[subsection]{Theorem}
\newtheorem{proposition}[subsection]{Proposition}
\newtheorem{lemma}[subsection]{Lemma}

\theoremstyle{definition}
\newtheorem{definition}[subsection]{Definition}
\newtheorem{example}[subsection]{Example}
\newtheorem{exercise}[subsection]{Exercise}
\newtheorem{situation}[subsection]{Situation}

\theoremstyle{remark}
\newtheorem{remark}[subsection]{Remark}
\newtheorem{remarks}[subsection]{Remarks}

\numberwithin{equation}{subsection}

% Macros
%
\def\lim{\mathop{\mathrm{lim}}\nolimits}
\def\colim{\mathop{\mathrm{colim}}\nolimits}
\def\Spec{\mathop{\mathrm{Spec}}}
\def\Hom{\mathop{\mathrm{Hom}}\nolimits}
\def\Ext{\mathop{\mathrm{Ext}}\nolimits}
\def\SheafHom{\mathop{\mathcal{H}\!\mathit{om}}\nolimits}
\def\SheafExt{\mathop{\mathcal{E}\!\mathit{xt}}\nolimits}
\def\Sch{\mathit{Sch}}
\def\Mor{\mathop{\mathrm{Mor}}\nolimits}
\def\Ob{\mathop{\mathrm{Ob}}\nolimits}
\def\Sh{\mathop{\mathit{Sh}}\nolimits}
\def\NL{\mathop{N\!L}\nolimits}
\def\CH{\mathop{\mathrm{CH}}\nolimits}
\def\proetale{{pro\text{-}\acute{e}tale}}
\def\etale{{\acute{e}tale}}
\def\QCoh{\mathit{QCoh}}
\def\Ker{\mathop{\mathrm{Ker}}}
\def\Im{\mathop{\mathrm{Im}}}
\def\Coker{\mathop{\mathrm{Coker}}}
\def\Coim{\mathop{\mathrm{Coim}}}

% Boxtimes
%
\DeclareMathSymbol{\boxtimes}{\mathbin}{AMSa}{"02}

%
% Macros for moduli stacks/spaces
%
\def\QCohstack{\mathcal{QC}\!\mathit{oh}}
\def\Cohstack{\mathcal{C}\!\mathit{oh}}
\def\Spacesstack{\mathcal{S}\!\mathit{paces}}
\def\Quotfunctor{\mathrm{Quot}}
\def\Hilbfunctor{\mathrm{Hilb}}
\def\Curvesstack{\mathcal{C}\!\mathit{urves}}
\def\Polarizedstack{\mathcal{P}\!\mathit{olarized}}
\def\Complexesstack{\mathcal{C}\!\mathit{omplexes}}
% \Pic is the operator that assigns to X its picard group, usage \Pic(X)
% \Picardstack_{X/B} denotes the Picard stack of X over B
% \Picardfunctor_{X/B} denotes the Picard functor of X over B
\def\Pic{\mathop{\mathrm{Pic}}\nolimits}
\def\Picardstack{\mathcal{P}\!\mathit{ic}}
\def\Picardfunctor{\mathrm{Pic}}
\def\Deformationcategory{\mathcal{D}\!\mathit{ef}}
\usepackage{cleveref}
\crefname{originallemma}{original lemma}{original lemmas}
\Crefname{originallemma}{Original lemma}{Original lemmas}
\makeatletter
\newlabel{lemma-abelian-variety-projective}{{0BFA}{source}{}{sourcepointer-1}{}}
\newlabel{lemma-abelian-variety-projective@cref}{{[originallemma][0][]0BFA}{[1][0][]source}{}{}{}}
\makeatother
\crefname{originallemma}{original lemma}{original lemmas}
\Crefname{originallemma}{Original lemma}{Original lemmas}
\makeatletter
\newlabel{more-morphisms-lemma-projective}{{0B45}{source}{}{sourcepointer-2}{}}
\newlabel{more-morphisms-lemma-projective@cref}{{[originallemma][0][]0B45}{[1][0][]source}{}{}{}}
\makeatother
\crefname{originallemma}{original lemma}{original lemmas}
\Crefname{originallemma}{Original lemma}{Original lemmas}
\makeatletter
\newlabel{properties-lemma-ample-tensor-globally-generated}{{0890}{source}{}{sourcepointer-3}{}}
\newlabel{properties-lemma-ample-tensor-globally-generated@cref}{{[originallemma][0][]0890}{[1][0][]source}{}{}{}}
\makeatother
\crefname{originallemma}{original lemma}{original lemmas}
\Crefname{originallemma}{Original lemma}{Original lemmas}
\makeatletter
\newlabel{lemma-pullbacks-by-n}{{0BFF}{source}{}{sourcepointer-4}{}}
\newlabel{lemma-pullbacks-by-n@cref}{{[originallemma][0][]0BFF}{[1][0][]source}{}{}{}}
\makeatother
\crefname{originallemma}{original lemma}{original lemmas}
\Crefname{originallemma}{Original lemma}{Original lemmas}
\makeatletter
\newlabel{varieties-lemma-ample-curve}{{0B5X}{source}{}{sourcepointer-5}{}}
\newlabel{varieties-lemma-ample-curve@cref}{{[originallemma][0][]0B5X}{[1][0][]source}{}{}{}}
\makeatother
\crefname{originallemma}{original lemma}{original lemmas}
\Crefname{originallemma}{Original lemma}{Original lemmas}
\makeatletter
\newlabel{varieties-lemma-ample-positive}{{0BEV}{source}{}{sourcepointer-6}{}}
\newlabel{varieties-lemma-ample-positive@cref}{{[originallemma][0][]0BEV}{[1][0][]source}{}{}{}}
\makeatother
\crefname{originallemma}{original lemma}{original lemmas}
\Crefname{originallemma}{Original lemma}{Original lemmas}
\makeatletter
\newlabel{coherent-lemma-characterize-finite}{{02OG}{source}{}{sourcepointer-7}{}}
\newlabel{coherent-lemma-characterize-finite@cref}{{[originallemma][0][]02OG}{[1][0][]source}{}{}{}}
\makeatother
\crefname{originallemma}{original lemma}{original lemmas}
\Crefname{originallemma}{Original lemma}{Original lemmas}
\makeatletter
\newlabel{lemma-abelian-variety-smooth}{{0BFC}{source}{}{sourcepointer-8}{}}
\newlabel{lemma-abelian-variety-smooth@cref}{{[originallemma][0][]0BFC}{[1][0][]source}{}{}{}}
\makeatother
\crefname{originallemma}{original lemma}{original lemmas}
\Crefname{originallemma}{Original lemma}{Original lemmas}
\makeatletter
\newlabel{algebra-lemma-CM-over-regular-flat}{{00R4}{source}{}{sourcepointer-9}{}}
\newlabel{algebra-lemma-CM-over-regular-flat@cref}{{[originallemma][0][]00R4}{[1][0][]source}{}{}{}}
\makeatother
\crefname{originallemma}{original lemma}{original lemmas}
\Crefname{originallemma}{Original lemma}{Original lemmas}
\makeatletter
\newlabel{varieties-lemma-smooth-regular}{{056S}{source}{}{sourcepointer-10}{}}
\newlabel{varieties-lemma-smooth-regular@cref}{{[originallemma][0][]056S}{[1][0][]source}{}{}{}}
\makeatother
\crefname{originallemma}{original lemma}{original lemmas}
\Crefname{originallemma}{Original lemma}{Original lemmas}
\makeatletter
\newlabel{algebra-lemma-regular-ring-CM}{{00NQ}{source}{}{sourcepointer-11}{}}
\newlabel{algebra-lemma-regular-ring-CM@cref}{{[originallemma][0][]00NQ}{[1][0][]source}{}{}{}}
\makeatother
\crefname{originallemma}{original lemma}{original lemmas}
\Crefname{originallemma}{Original lemma}{Original lemmas}
\makeatletter
\newlabel{morphisms-lemma-finite-flat}{{02KB}{source}{}{sourcepointer-12}{}}
\newlabel{morphisms-lemma-finite-flat@cref}{{[originallemma][0][]02KB}{[1][0][]source}{}{}{}}
\makeatother
\crefname{originallemma}{original lemma}{original lemmas}
\Crefname{originallemma}{Original lemma}{Original lemmas}
\makeatletter
\newlabel{varieties-lemma-degree-finite-morphism-in-terms-degrees}{{0BEX}{source}{}{sourcepointer-13}{}}
\newlabel{varieties-lemma-degree-finite-morphism-in-terms-degrees@cref}{{[originallemma][0][]0BEX}{[1][0][]source}{}{}{}}
\makeatother
\title{Multiplication by d on an abelian variety is finite locally free of degree d\textasciicircum{}\{2 dim\}}
\author{The Stacks Project Authors}
\date{stacks-project:groupoids.tex:local-snapshot-20261012; GNU FDL 1.2 or later}
\begin{document}
\maketitle
\setcounter{subsection}{1}
\noindent\textit{Editorial note.} All delivered mathematical text is the author's own, copied from the pinned Stacks Project chapter file groupoids.tex; no statement or proof step was rewritten, reordered or removed. The counted span is the author environment and its own proof environment. References inside the span are handed to the assembler as navigation stubs that carry the canonical tag of the cited tag only; whether each of them is load-bearing is recorded on the stub itself.\par\medskip
\noindent\textit{Allowed standard prerequisites (editorial).}\begin{itemize}
\item Lemmas lemma-abelian-variety-projective (0BFA) and more-morphisms-lemma-projective (0B45). cited at proof lines 1993-1995. Existence of the ample invertible module L.
\item Lemma properties-lemma-ample-tensor-globally-generated (0890) and lemma-pullbacks-by-n (0BFF). cited at proof lines 1998-2003. Ampleness of N and the identity [d]\textasciicircum{}*N = N\textasciicircum{}\{(x)d\textasciicircum{}2\}.
\item Varieties lemmas 0B5X (or 0BEV) and 02OG. cited at proof lines 2006-2014. Fibres of [d] are zero-dimensional, hence [d] finite.
\item Lemma lemma-abelian-variety-smooth (0BFC), algebra lemmas 00R4, 00NQ, varieties 056S, morphisms 02KB. cited at proof lines 2015-2024. Flatness of [d] and finite locally freeness.
\item Lemma varieties-lemma-degree-finite-morphism-in-terms-degrees (0BEX). cited at proof lines 2027-2031. The degree identity yielding d\textasciicircum{}\{2 dim(A)\}.
\end{itemize}\medskip
\ifdefined\hypertarget\hypertarget{sourcepointer-1}{}\fi
\noindent\textit{Source reference.} Original lemma 0BFA: sources/stacks-chow/tags.txt, source lines 14378--14378. Cited at proof lines 1993-1995 to supply the ample invertible module L; load-bearing.\par\medskip
\ifdefined\hypertarget\hypertarget{sourcepointer-2}{}\fi
\noindent\textit{Source reference.} Original lemma 0B45: sources/stacks-chow/tags.txt, source lines 13988--13988. Cited at proof lines 1993-1995 alongside 0BFA; load-bearing for the existence of L.\par\medskip
\ifdefined\hypertarget\hypertarget{sourcepointer-3}{}\fi
\noindent\textit{Source reference.} Original lemma 0890: sources/stacks-chow/tags.txt, source lines 10480--10480. Cited at proof lines 1998-2000 to see N = L (x) [-1]\textasciicircum{}*L is ample; load-bearing.\par\medskip
\ifdefined\hypertarget\hypertarget{sourcepointer-4}{}\fi
\noindent\textit{Source reference.} Original lemma 0BFF: sources/stacks-chow/tags.txt, source lines 14383--14383. Cited at proof lines 2001-2003 to get [d]\textasciicircum{}*N = N\textasciicircum{}\{(x)d\textasciicircum{}2\}; load-bearing for the degree formula.\par\medskip
\ifdefined\hypertarget\hypertarget{sourcepointer-5}{}\fi
\noindent\textit{Source reference.} Original lemma 0B5X: sources/stacks-chow/tags.txt, source lines 14050--14050. Cited at proof lines 2006-2010 to contradict O\_C ample of degree 0; load-bearing for finiteness.\par\medskip
\ifdefined\hypertarget\hypertarget{sourcepointer-6}{}\fi
\noindent\textit{Source reference.} Original lemma 0BEV: sources/stacks-chow/tags.txt, source lines 14363--14363. Cited parenthetically at proof line 2011 as an alternative to 0B5X; navigation for an alternative.\par\medskip
\ifdefined\hypertarget\hypertarget{sourcepointer-7}{}\fi
\noindent\textit{Source reference.} Original lemma 02OG: sources/stacks-chow/tags.txt, source lines 3515--3515. Cited at proof lines 2012-2014 to conclude [d] is finite; load-bearing.\par\medskip
\ifdefined\hypertarget\hypertarget{sourcepointer-8}{}\fi
\noindent\textit{Source reference.} Original lemma 0BFC: sources/stacks-chow/tags.txt, source lines 14380--14380. Cited at proof lines 2015-2017 to know A is smooth over k; load-bearing for flatness of [d].\par\medskip
\ifdefined\hypertarget\hypertarget{sourcepointer-9}{}\fi
\noindent\textit{Source reference.} Original lemma 00R4: sources/stacks-chow/tags.txt, source lines 990--990. Cited at proof lines 2017-2019 to see [d] is flat; load-bearing.\par\medskip
\ifdefined\hypertarget\hypertarget{sourcepointer-10}{}\fi
\noindent\textit{Source reference.} Original lemma 056S: sources/stacks-chow/tags.txt, source lines 6720--6720. Cited at proof lines 2019-2021 in the regularity step of the flatness argument; load-bearing.\par\medskip
\ifdefined\hypertarget\hypertarget{sourcepointer-11}{}\fi
\noindent\textit{Source reference.} Original lemma 00NQ: sources/stacks-chow/tags.txt, source lines 868--868. Cited at proof lines 2020-2021 in the regularity step; load-bearing.\par\medskip
\ifdefined\hypertarget\hypertarget{sourcepointer-12}{}\fi
\noindent\textit{Source reference.} Original lemma 02KB: sources/stacks-chow/tags.txt, source lines 3366--3366. Cited at proof lines 2022-2024 to conclude [d] is finite locally free; load-bearing.\par\medskip
\ifdefined\hypertarget\hypertarget{sourcepointer-13}{}\fi
\noindent\textit{Source reference.} Original lemma 0BEX: sources/stacks-chow/tags.txt, source lines 14365--14365. Cited at proof lines 2027-2031 for the degree identity deg(N\textasciicircum{}\{(x)d\textasciicircum{}2\}) = deg([d]) deg(N); load-bearing for the final formula.\par\medskip
\begin{lemma}
\label{lemma-degree-multiplication-by-d}
Let $k$ be a field. Let $A$ be an abelian variety over $k$.
Let $[d] : A \to A$ be the multiplication by $d$.
Then $[d]$ is finite locally free of degree $d^{2\dim(A)}$.
\end{lemma}

\begin{proof}
By Lemma \ref{lemma-abelian-variety-projective}
(and More on Morphisms, Lemma \ref{more-morphisms-lemma-projective})
we see that $A$ has an ample invertible module $\mathcal{L}$.
Since $[-1] : A \to A$ is an automorphism, we see that
$[-1]^*\mathcal{L}$ is an ample invertible $\mathcal{O}_A$-module
as well. Thus $\mathcal{N} = \mathcal{L} \otimes [-1]^*\mathcal{L}$
is ample, see
Properties, Lemma \ref{properties-lemma-ample-tensor-globally-generated}.
Since $\mathcal{N} \cong [-1]^*\mathcal{N}$ we see that
$[d]^*\mathcal{N} \cong \mathcal{N}^{\otimes d^2}$ by
Lemma \ref{lemma-pullbacks-by-n}.

\medskip\noindent
To get a contradiction, let $C$ be a proper curve contained in a
fibre of $[d]$. Then $\mathcal{N}^{\otimes d^2}|_C \cong \mathcal{O}_C$
is an ample invertible $\mathcal{O}_C$-module of degree $0$ which
contradicts Varieties, Lemma \ref{varieties-lemma-ample-curve} for example.
(You can also use Varieties, Lemma \ref{varieties-lemma-ample-positive}.)
Thus every fibre of $[d]$ has dimension $0$ and hence $[d]$ is finite
for example by Cohomology of Schemes, Lemma
\ref{coherent-lemma-characterize-finite}.
Moreover, since $A$ is smooth over $k$ by
Lemma \ref{lemma-abelian-variety-smooth}
we see that $[d] : A \to A$ is flat by
Algebra, Lemma \ref{algebra-lemma-CM-over-regular-flat}
(we also use that schemes smooth over fields are regular and that
regular rings are Cohen-Macaulay, see
Varieties, Lemma \ref{varieties-lemma-smooth-regular} and
Algebra, Lemma \ref{algebra-lemma-regular-ring-CM}).
Thus $[d]$ is finite flat hence finite locally free by
Morphisms, Lemma \ref{morphisms-lemma-finite-flat}.

\medskip\noindent
Finally, we come to the formula for the degree. By
Varieties, Lemma \ref{varieties-lemma-degree-finite-morphism-in-terms-degrees}
we see that
$$
\deg_{\mathcal{N}^{\otimes d^2}}(A) = \deg([d]) \deg_\mathcal{N}(A)
$$
Since the degree of $A$ with respect to
$\mathcal{N}^{\otimes d^2}$, respectively $\mathcal{N}$
is the coefficient of $n^{\dim(A)}$ in the polynomial
$$
n \longmapsto \chi(A, \mathcal{N}^{\otimes nd^2}),\quad
\text{respectively}\quad n \longmapsto \chi(A, \mathcal{N}^{\otimes n})
$$
we see that $\deg([d]) = d^{2 \dim(A)}$.
\end{proof}

\end{document}
